\chapter{Main Results}


\section{Summary of equilibrium notions}

\subsection*{Table 1: Walrasian and Financial Market equilibria}

\renewcommand{\arraystretch}{1.3}
\setlength{\tabcolsep}{5pt}
\small
\noindent\begin{tabularx}{\textwidth}{|>{\raggedright\arraybackslash}p{3cm}|>{\raggedright\arraybackslash}p{2.8cm}|>{\raggedright\arraybackslash}X|>{\raggedright\arraybackslash}p{3.4cm}|}
\hline
\textbf{Equilibrium} & \textbf{Variables} & \textbf{Object space} & \textbf{Efficiency} \\
\hline
Abstract Exchange CE & $(x, p)$ & $\mathbb{R}_+^{LI} \times \mathbb{R}_{++}^L$ & 1st \& 2nd WT \\
\hline
Exchange w/ Production & $(x, y, p)$ & $\mathbb{R}_+^{LI} \times Y \times \mathbb{R}_{++}^L$ & 1st \& 2nd WT \\
\hline
Arrow--Debreu (exchange) & $(x, \phi)$ & $\mathbb{R}_+^{L(S+1)I} \times \mathbb{R}_{++}^{L(S+1)}$ & 1st \& 2nd WT \\
\hline
Arrow--Debreu w/ Production & $(x, y, \phi)$ & $\mathbb{R}_+^{L(S+1)I} \times Y \times \mathbb{R}_{++}^{L(S+1)}$ & 1st \& 2nd WT \\
\hline
Financial Market Eq.\ (FME) & $(x, z, p, q)$ & $\mathbb{R}_+^{L(S+1)I} \times \mathbb{R}^{JI} \times \mathbb{R}_{++}^{L(S+1)} \times \mathbb{R}^J$ & Constr.\ PO if $L=1$; fails if $L>1$ (G--P) \\
\hline
FME with No-Arbitrage & $(x, \pi, p)$ & $\mathbb{R}_+^{L(S+1)I} \times \mathbb{R}_{++}^S \times \mathbb{R}_{++}^{L(S+1)}$ & = FME (reformulation) \\
\hline
FME + Production + Rational Conj. & $(x, k, z, \theta, p, q, Q)$ & FME dim.\ $\times\ K^H \times \mathbb{R}_+^{HI} \times (\mathbb{R}_+^K)^H$ & Generically not constr.\ PO \\
\hline
Representative Agent FME & $(x^R, z^R, p, q)$ & FME w/ $I=1$; $x^R = \sum_i \omega^i$, $z^R = 0$ & Inherits FME; autarchic \\
\hline
\end{tabularx}


\subsection*{Table 2: Dynamic, sunspot, and asymmetric-information equilibria}

\renewcommand{\arraystretch}{1.3}
\setlength{\tabcolsep}{5pt}
\small
\noindent\begin{tabularx}{\textwidth}{|>{\raggedright\arraybackslash}p{3cm}|>{\raggedright\arraybackslash}p{2.8cm}|>{\raggedright\arraybackslash}X|>{\raggedright\arraybackslash}p{3.4cm}|}
\hline
\textbf{Equilibrium} & \textbf{Variables} & \textbf{Object space} & \textbf{Efficiency} \\
\hline

Sunspot Equilibrium & $(x, z, p, q)$ on $S \times \Sigma$ & FME dim., indexed by $\sigma \in \Sigma$ & Generically not PO \\
\hline

Sequential FME & $\{(x^i, z^i)_i, q\}$ per node $s^t$ & Indexed by event tree $\mathcal{N}$ & Static FME node-by-node \\
\hline

Sequential Arrow--Debreu & $(x^*, p^*)$ per node $s^t$ & $\mathbb{R}_+^{LI \cdot |\mathcal{N}|} \times \mathbb{R}_{++}^{L \cdot |\mathcal{N}|}$ & 1st \& 2nd WT \\
\hline

Rational Exp.\ Eq.\ (REE) & $\phi: \Sigma^I \to \mathbb{R}_+^J$ + alloc. & Price-map space $\{\phi: \Sigma^I \to \mathbb{R}_+^J\}$ & Info aggregation; not PO \\
\hline

P--T (Moral Hazard) & $(x, e, q(\cdot))$ & $\mathbb{R}_+^S \times E \times \{q: E \to \mathbb{R}_+^S\}$ + IC & Incentive-constr.\ PO \\
\hline

P--T (Adverse Selection) & $(x^e)_{e \in E}, q(\cdot)$ & $\mathbb{R}_+^{SE} \times \{q: E \to \mathbb{R}_+^S\}$ + truth-telling & Generically \emph{not} IC-PO \\
\hline
\end{tabularx}

\section{Concepts}

\subsection{Boundary condition}
The usual boundary condition used in the existence proof is:
\[
p^n \rightarrow p \text{ with some } p_l = 0 \Rightarrow \max \{z_1(p^n, w), z_2(p^n, w), 
dots, z_L(p^n, w) \} \rightarrow \infty.
\]

\subsection{Aggregate Excess Demand Function Properties}
Let $z(p, w): \mathbb{R}_{++}^{L} \times \mathbb{R}_{++}^{LI} \rightarrow \mathbb{R}_{+}^{L}$ be the excess demand function. 
It has the following properties ("SWBBH")
\begin{enumerate}
    \item smoothness: $z(p, w)$ is $C^1$
    \item Walras Law: $p z(p, w) = 0$.
    \item bounded from below: $z(p, w) \geq -s $
    \item boundary condition: $p^n \rightarrow p$ with some $p_l = 0$ then,$\max \{z_k(p, w) \} \rightarrow \infty$. 
    \item homogeneity of degree $0$: $z(\lambda p, \lambda w) = \lambda^0 z(p, w) = z(p, w)$.
\end{enumerate}

\section{Propositions}

\subsection{Sunspots}

Intuitively, a sunspot equilibrium means that we can have different equilibrium outcomes even when economic fundamentals are the same.

\begin{tcolorbox}
$\exists$ non-trivial sunspot equilibrium 
$\Rightarrow$
 equilibrium indeterminacy or multiplicity.
\end{tcolorbox}

\paragraph{Cass and Shell.}
When markets are complete, sunspot equilibria do not exist. That is, if
\[
\omega_{s,e}=\omega_{s,e'}=\omega_s,
\]
then
\[
x_{s,e}=x_{s,e'}.
\]
But when markets are incomplete, sunspot equilibria may exist.

\paragraph{Example.}
Since there is only one bond paying one unit of the numeraire in every state,
the asset payoff vector is \(A=(1,\dots,1)'\). Hence, the span condition implies
\[
p_s(x_s^i-\omega_s^i)=z^i
\qquad \text{for every state }s.
\]
After splitting state \(1\) into two extrinsic states \(\alpha,\beta\), we have
\[
p_\alpha(x_\alpha^i-\omega_1^i)
=
p_\beta(x_\beta^i-\omega_1^i)
=
z^i.
\]
Thus, the same bond position must finance net trades in both extrinsic states.
But this condition only restricts the value of net trades, not the commodity bundle
itself. Therefore it does not force \(x_\alpha^i=x_\beta^i\).

If equilibrium spot prices differ across the two extrinsic states, \(p_\alpha\neq p_\beta\),
then the same numeraire payoff can support different consumption bundles. Hence,
with incomplete markets, it is possible that
\[
x_\alpha\neq x_\beta
\]
even though \(\omega_\alpha=\omega_\beta=\omega_1\).
This is the sunspot logic.


\subsection{FWT under Asymmetric Information}
Rationality of conjectures is critical for the First welfare theorem in the simple insurance economy. 